% SEGMENT OLP-0569-S01
\documentclass[../../../include/open-logic-section]{subfiles}

\begin{document}

\olfileid{sth}{replacement}{limofsize}
\olsection{அளவு-வரம்பு}

மாற்றீட்டுக்கு ஓர் ``உள்ளார்ந்த'' நியாயப்படுத்தலை வழங்கும் மிகப் பொதுவான
முயற்சி பின்வரும் கருத்தின் வழியே வருவதாக இருக்கலாம்:
\begin{enumerate}
	\item[] \limofsize. பொருள்கள் மிக அதிகமாக இல்லாதவரை, எந்தப் பொருள்களும்
	ஒரு கணத்தை அமைக்கின்றன.
\end{enumerate}
இக்கொள்கை மாற்றீட்டை உடனடியாக உறுதிசெய்யும். ஏனெனில் மாற்றீட்டால் உருவான
எந்தக் கணமும், அது எந்தக் கணத்திலிருந்து உருவானதோ அதைவிடப் பெரியதாக இருக்க
முடியாது. துல்லியமாகக் கூறினால்: மாற்றீட்டைப் பயன்படுத்தி
$\funimage{\tau}{A} = \Setabs{\tau(x)}{x \in A}$ என்ற கணத்தை நீங்கள்
உருவாக்குகிறீர்கள் என்க; அப்போது $\cardle{\funimage{\tau}{A}}{A}$; எனவே
$A$ இன் !!{element}s ஒரு கணத்தை அமைக்க முடியாத அளவுக்கு அதிகமில்லை என்றால்,
அவற்றின் பிம்பங்களும் $\funimage{\tau}{A}$ ஐ அமைக்க முடியாத அளவுக்கு
அதிகமில்லை.

மாற்றீட்டை நியாயப்படுத்த \limofsize{} ஐப் பயன்படுத்துவதிலுள்ள வெளிப்படையான
இடர்ப்பாடு, \limofsize{} போன்ற எந்தக் கொள்கையையும் நாம் \emph{இன்னும்}
வகுக்கவில்லை என்பதாகும். மேலும்,
\crefrange{sth:story::chap}{sth:z::chap} இல் கணத்தின்
திரள்-படிநிலைமுறைக் கருத்தாக்கத்தைப் பற்றிய நமது கதையைச் சொன்னபோது,
\limofsize{} ஐ நோக்கி எதுவும் \emph{சுட்டிக்காட்டவே} இல்லை. இதனால்தான்
பூலோஸ் ஒருகட்டத்தில் ``கணக் கோட்பாட்டுக்குப் `பின்னால்' குறைந்தபட்சம் இரண்டு
சிந்தனைகள் உள்ளன என்று ஒருவர் முடிவுசெய்யலாம்'' என்று எழுதினார்
\citeyearpar[p.~19]{Boolos1989}. ஒருபுறம், கணத்தின் திரள்-படிநிலைமுறைக்
கருத்தாக்கத்தைச் சூழ்ந்த கருத்துகள்~$\Z$ ஐ உறுதிசெய்வதாக உள்ளன. மறுபுறம்,
\limofsize{} மாற்றீட்டை உறுதிசெய்வதாக உள்ளது.

ஆனால் இங்குள்ள சிக்கல், இதுவரை \limofsize{} பற்றி நாம் \emph{மௌனமாக}
இருந்தோம் என்பது மட்டுமன்று. மாறாக, இப்போது வகுத்துள்ளபடி \limofsize{}
ஆனது கணத்தின் திரள்-படிநிலைமுறைக் கருத்துடன் மிக மோசமாகப் பொருந்துவதாகத்
தோன்றுகிறது. ஏனெனில் கணங்கள் \emph{படிநிலைகளில்} உருவாகின்றன என்ற கருத்தை
அது சிறிதும் குறிப்பிடவில்லை.
% SOURCE CORRECTION TA-STH-018: remove the stray “it” from “the issue it is not just.”

இவ்விரு ``சிந்தனைகளின்'' வரலாற்றைக் கருத்தில் கொண்டால் இது உண்மையில்
வியப்புக்குரியதன்று (அதாவது, கணத்தின் திரள்-படிநிலைமுறைக் கருத்தாக்கமும்
\limofsize{} உம்). இவ்விரு ``சிந்தனைகளும்'' இறுதியில் கணக்கோட்பாட்டு
முரணுரைகளைத் தடுக்கும் முற்றிலும் வேறுபட்ட இரு திட்டங்களாக அமைகின்றன.
கணத்தின் திரள்-படிநிலைமுறைக் கருத்து, ஏறத்தாழப் பின்வருமாறு கூறி ரஸ்ஸலின்
முரணுரையைத் தடுக்கிறது: \emph{ரஸ்ஸல் கணம் எந்தப் படிநிலையிலும்
``உருவாகாது'' என்பதால், அது இருக்கும் என்று நாம் ஒருபோதும்
எதிர்பார்த்திருக்கக் கூடாது}. இதற்கு மாறாக, \limofsize{} பின்வருமாறு கூறி
ரஸ்ஸல் கணத்தை விலக்குவதற்கானது: \emph{ரஸ்ஸல் கணம் மிகவும் பெரியதாக
இருக்கும் என்பதால், அது இருக்கும் என்று நாம் ஒருபோதும் எதிர்பார்த்திருக்கக்
கூடாது}.

இவ்வாறு கூறினால், வெளிப்படையாகப் பேசுவோம்: முரணுரைகளுக்கான பதிலாகக்
கருதும்போது, \limofsize{} க்கு \emph{மிகவும்} அதிகமான நியாயப்படுத்தல்
தேவை. எடுத்துக்காட்டாக, ரஸ்ஸலின் முரணுரையின் இந்த வடிவத்தைக் கருதுக:
\emph{தம்மைத்தாமே முகராத குட்டைநாய்களை, அவற்றை மட்டுமே, முகரும் குட்டைநாய்
எதுவும் இல்லை} (\olref[sth][story][rus]{sec} ஐப் பார்க்கவும்). ``அத்தகைய
குட்டைநாய் ஏன் இல்லை?'' என்று நீங்கள் கேட்டால், அது மிக அதிகமான
குட்டைநாய்களை முகர வேண்டியிருக்கும் என்று கூறுவது நல்ல விடையன்று. அப்படியானால்,
ரஸ்ஸல் கணம் இருக்காததற்கான நல்ல உள்ளுணர்வு விளக்கம், அது இருக்க முடியாத
அளவுக்கு ``மிகப் பெரியதாக'' இருக்க வேண்டும் என்பது ஏன்?

சுருக்கமாக, \limofsize{} க்கான ``உள்ளுணர்வு'' ஊக்கம் குறித்து நீங்கள்
சற்றுக் குழம்பினால் அது மன்னிக்கத்தக்கதே.

\end{document}
